% Required preamble packages for this chapter:
%   amsmath, booktabs, graphicx, multirow, glossaries
% Acronyms assumed defined: ga, sade, aiwpso, gwo, bo, cnn, nas, hpo
% New BibTeX keys used here (see the end-of-chapter note in the progress log):
%   dong2020nasbench201, lecun1998mnist, krizhevsky2009cifar, goldberg1989genetic,
%   qin2009sade, qin2006adaptive, mirjalili2014gwo, snoek2012bo, deb2002nsga2,
%   kukkonen2005gde3, daulton2020qehvi, zitzler1999hv, ishibuchi2015igdplus,
%   demsar2006statistical, van2023mealpy

\chapter{Experimental Evaluation}
\label{cha:experimental_evaluation}

This chapter reports the experimental evaluation of the approach outlined in this work, which applies evolutionary optimization to different parts of the machine learning pipeline. The evaluation is organized as three studies that share the same set of optimization algorithms, the same seed and budget protocol, and the same reporting conventions, so that the only quantity that changes between them is the problem formulation. This design makes it possible to ask not only \emph{which algorithm is best}, but also \emph{whether an algorithm ranking established on one search space transfers to another}.

Section~\ref{sec:experimental_environment} describes the computational environment, the algorithms compared across the studies, their control parameters, and the conventions shared by all experiments.

Section~\ref{sec:study_1} presents Study~1, which addresses \acrfull{hpo} of a fixed \acrshort{cnn} template over a categorical space of seven hyperparameters, on MNIST~\cite{lecun1998mnist} and CIFAR-10~\cite{krizhevsky2009cifar}. Every candidate is decoded into a network that is actually built and trained, so each fitness evaluation carries a real training cost.

Section~\ref{sec:study_2} presents Study~2, which moves from hyperparameters to topology and addresses \acrfull{nas} on the tabular benchmark NAS-Bench-201~\cite{dong2020nasbench201}. Because every architecture in this benchmark is pre-trained, each evaluation reduces to a lookup and the global optimum is known by construction. The study is organized in three parts: an accuracy-only single-objective formulation, a cost-aware formulation that scalarizes accuracy against computational cost, and a native multi-objective formulation that searches directly for the Pareto front.

Section~\ref{sec:study_3} outlines Study~3, which targets the third role of evolutionary methods inside the learning phase, the optimization of the network weights themselves, through a hybrid scheme that combines a population-based global search with gradient-based local refinement. This study is still under definition and is presented here as a plan rather than as a set of results.

Section~\ref{sec:experimental_summary} consolidates the three studies and discusses the findings that emerge from reading them together.


\section{Experimental Environment}
\label{sec:experimental_environment}

All experiments are executed in a Python~3 environment. Study~1 uses TensorFlow/Keras as the deep learning framework, since it builds and trains a network for every fitness evaluation. Study~2 does not train any network during the search and instead queries the NAS-Bench-201 tabular benchmark. Table~\ref{tab:hardware} summarizes the hardware and software environment.

\begin{table}[ht]
    \centering
    \caption{Hardware and software environment used in the experiments.}
    \label{tab:hardware}
    \begin{tabular}{ll}
        \toprule
        \textbf{Component}    & \textbf{Specification} \\
        \midrule
        CPU                   & AMD Ryzen 7 7700X 8-Core Processor \\
        RAM                   & 32.0 GB \\
        GPU                   & NVIDIA GeForce RTX 5060 Ti \\
        Operating system      & Ubuntu 22.04 LTS \\
        Python version        & 3.11 \\
        Deep learning         & TensorFlow / Keras \textit{---} \\ % TODO: confirm the exact TF/Keras version used in Study 1
        Metaheuristics        & \texttt{mealpy}~\cite{van2023mealpy} \\
        Bayesian optimization & \texttt{BoTorch} 0.16.1 \\
        Multi-objective       & \texttt{pymoo} 0.6.1.6 \\
        Benchmark             & NAS-Bench-201 API 1.8 \\
        \bottomrule
    \end{tabular}
\end{table}

\subsection{Optimization Algorithms}
\label{subsec:algorithms}

The set of algorithms compared is the same across studies. It was chosen to span the main algorithmic families used in evolutionary machine learning rather than several variants of a single family, so that the comparison probes whether an algorithm's family, and not merely its tuning, drives its relative performance. Four are evolutionary or swarm metaheuristics and one is a non-evolutionary baseline included as a strong reference:

\begin{enumerate}
    \item \acrfull{ga}~\cite{goldberg1989genetic}, which evolves fixed-length encodings through tournament selection, uniform crossover and flip mutation;
    \item \acrfull{sade}~\cite{qin2009sade}, which self-adapts its scaling factor, its crossover rate and even its mutation strategy from the success history of previous trials, exposing no user-facing knobs beyond the population size;
    \item \acrfull{aiwpso}~\cite{qin2006adaptive}, which moves a swarm of particles toward the personal and global best positions with an inertia weight adapted from each particle's fitness rank;
    \item \acrfull{gwo}~\cite{mirjalili2014gwo}, which simulates the grey wolf leadership hierarchy and hunting behavior, and whose coefficients are deterministic functions of the iteration index, leaving no free parameters;
    \item \acrfull{bo}~\cite{snoek2012bo}, the non-evolutionary baseline, which builds a Gaussian-process surrogate of the objective and selects each query by maximizing an Expected Improvement acquisition function, never re-evaluating a point.
\end{enumerate}

The multi-objective part of Study~2 (Section~\ref{subsec:study2_multiobjective}) additionally requires optimizers that return a set of non-dominated solutions rather than a single one. Three native Pareto optimizers are used there: Non-dominated Sorting Genetic Algorithm~II (NSGA-II)~\cite{deb2002nsga2}, which combines fast non-dominated sorting with a crowding-distance operator; Generalized Differential Evolution~3 (GDE3)~\cite{kukkonen2005gde3}, which extends differential evolution with a one-to-one Pareto selection and the same diversity-preserving truncation; and the Bayesian q-Expected Hypervolume Improvement (qEHVI)~\cite{daulton2020qehvi}, which fits a Gaussian-process surrogate and proposes, at each step, the candidate expected to enlarge the dominated hypervolume the most.

This work deliberately favors formulations that \emph{adapt their own control parameters during the optimization process}, since a method that must itself be tuned is of limited practical use in exactly the setting where automated search is needed. \acrshort{sade} is the clearest example: the scaling factor, the crossover rate and the mutation strategy are continuously updated from past trials. \acrshort{gwo} reaches the same practical outcome by a different route, exposing no control parameters at all.

Every configurable parameter is kept at its literature-recommended default rather than being tuned per problem, so as not to add a second, confounding layer of optimization on top of the one under study. Table~\ref{tab:algorithm_params} lists the resulting settings.

\begin{table}[ht]
    \centering
    \caption{Control parameters of the algorithms under comparison. All population-based methods use a population of $20$ individuals. Values are the library or literature defaults and are kept fixed across every repetition and every study.}
    \label{tab:algorithm_params}
    \begin{tabular}{ll}
        \toprule
        \textbf{Algorithm} & \textbf{Control parameters} \\
        \midrule
        \acrshort{ga}      & tournament selection ($k = 0.2$), uniform crossover ($p_c = 0.9$), \\
                           & flip mutation ($p_m = 0.05$) \\
        \acrshort{sade}    & none (self-adaptive $F$, $CR$ and mutation strategy) \\
        \acrshort{aiwpso}  & $c_1 = c_2 = 2.05$, inertia-weight rank exponent $\alpha = 0.4$ \\
        \acrshort{gwo}     & none (coefficients are deterministic functions of the iteration index) \\
        \acrshort{bo}      & Gaussian-process surrogate, Expected Improvement with $\xi = 0.01$ \\
        \midrule
        NSGA-II            & fast non-dominated sorting, crowding-distance truncation \\
        GDE3               & one-to-one Pareto selection on a differential-evolution population \\
        qEHVI              & Gaussian-process surrogate, q-Expected Hypervolume Improvement \\
        \bottomrule
    \end{tabular}
\end{table}

The four evolutionary and swarm methods are implemented with the \texttt{mealpy} library~\cite{van2023mealpy}, \acrshort{bo} with a Gaussian-process implementation, and the native multi-objective optimizers with \texttt{pymoo} (NSGA-II, GDE3) and \texttt{BoTorch} (qEHVI).

\subsection{Common Conventions}
\label{subsec:common_conventions}

Three conventions hold across all studies and are stated here once.

First, \emph{fitness is always expressed as a quantity to minimize}. Since the target metric is an accuracy, the fitness of a candidate $\mathbf{x}$ is always of the form $1 - \mathrm{acc}(\mathbf{x})$, so that lower values are better and the convergence plots read downward.

Second, \emph{accuracy is used throughout} as the fitness signal and as the reported metric. This choice relies on all datasets involved having balanced classes by construction, including the 120-class ImageNet16-120 subset bundled with NAS-Bench-201.

Third, \emph{every run is seeded and repeated}. A seed fixes every source of stochasticity of the run it controls: the optimizer's own random operators in all studies, plus the network's weight initialization and batch order in Study~1, where a network is actually trained during the search. Study~2 candidates are instead scored by an already-trained benchmark, so the only stochasticity is the optimizer's. For every setting the aggregated convergence behavior, the distribution of the final result across repetitions, and condensed summary statistics are reported, rather than a single best run.


\section{Study 1 --- Hyperparameter Optimization of a CNN Template}
\label{sec:study_1}

The first study addresses the question \emph{``which optimization algorithm is most effective at navigating a categorical \acrshort{cnn} hyperparameter space under a fixed and genuinely expensive evaluation budget?''}. The architecture is fixed by a template and the optimizer chooses a subset of its hyperparameters. Every candidate is decoded into a real network that is built and trained, so the cost of the search is dominated by the training of the candidates rather than by the optimizer.

\subsection{Datasets}
\label{subsec:datasets}

Two image-classification datasets are used, both balanced and computationally tractable, and both standard enough that the results are comparable with the literature.

\begin{enumerate}
    \item \textbf{MNIST}~\cite{lecun1998mnist}: grayscale images of the handwritten digits 0--9, with $60{,}000$ images in the training set and $10{,}000$ in the test set, evenly distributed across the ten classes.
    \item \textbf{CIFAR-10}~\cite{krizhevsky2009cifar}: RGB images of ten object categories, with $50{,}000$ images in the training set and $10{,}000$ in the test set, again evenly distributed.
\end{enumerate}

The same \acrshort{cnn} template (Section~\ref{subsec:architecture_design}), the same hyperparameter search space (Section~\ref{subsec:search_space}) and the same optimization pipeline (Section~\ref{sec:optimisation_pipeline}) are used on both datasets. The only differences are the input shape and the resulting number of model parameters. Table~\ref{tab:datasets} summarizes this information.

\begin{table}[h]
    \centering
    \caption{Datasets used in Study~1.}
    \label{tab:datasets}
    \begin{tabular}{lcccc}
        \toprule
        \textbf{Dataset} & \textbf{Labels} & \textbf{Input shape} & \textbf{Train / Test size} & \textbf{\# HP} \\
        \midrule
        MNIST    & 10 & $28\times28\times1$ & $60{,}000$ / $10{,}000$ & 7 \\
        CIFAR-10 & 10 & $32\times32\times3$ & $50{,}000$ / $10{,}000$ & 7 \\
        \bottomrule
    \end{tabular}
\end{table}

The two datasets play different roles. MNIST is included because it is the standard entry point for this kind of comparison, but it is close to saturation for any reasonable \acrshort{cnn}, and this turns out to matter for the interpretation of the results. CIFAR-10 leaves genuine room for improvement under the same budget, and is therefore the benchmark that actually discriminates between optimizers.

\subsection{Architecture Design}
\label{subsec:architecture_design}

Allowing the optimizer to assemble an arbitrary graph of layers would lead to a combinatorial explosion of the search space, because each candidate architecture would be defined not only by the type and hyperparameters of every layer, but also by the connectivity between them. Exploring such a space meaningfully requires a fitness budget that exceeds the computational resources available for this work: evaluating a decoded architecture requires training it, and even a reduced training budget per candidate becomes impractical as the number of candidates grows.

For this reason a fixed \emph{template architecture} is adopted, and the optimizer is restricted to choosing a subset of its hyperparameters. The template is a sequential \acrshort{cnn} made up of $X$ convolutional \textit{blocks} followed by a classifier head, as illustrated in Figure~\ref{fig:buildTemplateCNN_architecture_2}.

\begin{figure}[ht]
    \centering
    \includegraphics[width=\textwidth]{Chapters/img/buildTemplateCNN_architecture_2.pdf}
    \caption{Example of a decoded \acrshort{cnn} architecture produced by the template.}
    \label{fig:buildTemplateCNN_architecture_2}
\end{figure}

Each convolutional block is built with \emph{two stacked convolutional layers} of the same width, a design pattern that recurs in many well-established architectures for image classification such as VGG, ResNet and DenseNet. Concretely, each block is composed of the following sequence:

\begin{enumerate}
    \item A $k \times k$ convolution with \texttt{same} padding, followed by a batch-normalization layer and a non-linear activation;
    \item A second $k \times k$ convolution of the same width, followed again by batch normalization and activation;
    \item A $2 \times 2$ spatial pooling operation, either max-pooling or average-pooling, that halves the feature-map resolution.
\end{enumerate}

The number of filters in the $i$-th block grows geometrically with depth, $f_i = f_0 \cdot 2^{i}$, where $f_0$ is the base number of filters. This progressive widening compensates for the reduction in spatial resolution introduced by pooling and is a standard choice across \acrshort{cnn} architectures.

After the last convolutional block, the classifier head flattens the feature maps, applies a dense layer with a configurable number of units, a non-linear activation and a dropout, and finishes with a softmax layer with one output per class.

Two design choices are kept \emph{fixed} on methodological grounds and are therefore not exposed to the optimizer:

\begin{enumerate}
    \item Batch normalization is always applied after every convolution, since its regularizing and optimization benefits are well documented and do not interact with the hyperparameters being searched;
    \item The geometric progression of the number of filters is treated as a design principle rather than a free dimension, so the only degree of freedom on this axis is the base number of filters $f_0$.
\end{enumerate}

\subsection{Search Space}
\label{subsec:search_space}

The search space of Study~1 exposes seven hyperparameters with discrete values and is summarized in Table~\ref{tab:search_space_study1}. Each candidate solution is encoded as an integer vector of length seven, in which every coordinate is the index of the chosen option in its corresponding list. This categorical encoding keeps the search space small and discrete, which is appropriate for the short fitness budget imposed by the cost of training each candidate.

\begin{table}[ht]
    \centering
    \caption{Search space of Study~1. Each candidate solution is encoded as a length-7 integer vector, giving a cardinality of $3^{4} \cdot 2^{3} = 648$ configurations.}
    \label{tab:search_space_study1}
    \begin{tabular}{clcc}
        \toprule
        \textbf{Index} & \textbf{Hyperparameter} & \textbf{Possible values} & \textbf{\# options} \\
        \midrule
        0 & Batch size              & $\{32, 64, 128\}$                & 3 \\
        1 & Base filters $f_0$      & $\{32, 64, 128\}$                & 3 \\
        2 & Number of blocks $X$    & $\{2, 3\}$                       & 2 \\
        3 & Dense units             & $\{128, 256, 512\}$              & 3 \\
        4 & Optimizer               & $\{\text{AdamW}, \text{SGD}, \text{RMSprop}\}$ & 3 \\
        5 & Pooling type            & $\{\text{max}, \text{average}\}$ & 2 \\
        6 & Kernel size $k$         & $\{3, 5\}$                       & 2 \\
        \midrule
        \multicolumn{3}{l}{\textbf{Cardinality of the search space}} & $648$ \\
        \bottomrule
    \end{tabular}
\end{table}

Everything that is not in Table~\ref{tab:search_space_study1} is held fixed across every run and every algorithm, as listed in Table~\ref{tab:fixed_hyperparams}. Holding these values constant is what makes the comparison between optimizers meaningful: any difference in the reported fitness is attributable to the seven searched coordinates and to the optimizer that chose them.

\begin{table}[ht]
    \centering
    \caption{Hyperparameters held fixed across all runs of Study~1.}
    \label{tab:fixed_hyperparams}
    \begin{tabular}{ll}
        \toprule
        \textbf{Hyperparameter} & \textbf{Fixed value} \\
        \midrule
        Activation function          & ReLU \\
        Batch normalization          & enabled (after every convolution) \\
        Dropout rate                 & $0.5$ \\
        Learning rate                & $10^{-3}$ \\
        Loss function                & categorical cross-entropy \\
        \bottomrule
    \end{tabular}
\end{table}

\subsection{Optimization Pipeline}
\label{sec:optimisation_pipeline}

The optimization pipeline drives a metaheuristic through the search space defined in Section~\ref{subsec:search_space} and returns the best encoding it can find. The overall workflow is depicted in Figure~\ref{fig:hpo_workflow}.

\begin{figure}[ht]
    \centering
    \includegraphics[width=\textwidth]{Chapters/img/hpo_workflow.pdf}
    \caption{Workflow of the Study~1 optimization pipeline. Each candidate encoding is decoded into a \acrshort{cnn} configuration, built, trained under a reduced budget, and scored; the best configuration is then rebuilt and fully retrained.}
    \label{fig:hpo_workflow}
\end{figure}

The pipeline operates as a closed loop. The metaheuristic proposes an integer-encoded candidate; the encoding is decoded into a \acrshort{cnn} hyperparameter configuration; the corresponding network is built and trained for a single epoch on the selected dataset; and the resulting validation accuracy is returned to the algorithm as the fitness signal that guides the next proposal. The fitness of a candidate $\mathbf{x}$ is

\begin{equation}
    \mathrm{fit}(\mathbf{x}) \;=\; 1 - \mathrm{acc}_{\mathrm{val}}(\mathbf{x}),
    \label{eq:fitness_study1}
\end{equation}

where $\mathrm{acc}_{\mathrm{val}}$ denotes the validation accuracy of the \acrshort{cnn} trained with the hyperparameters $\mathbf{x}$. Lower fitness is therefore better.

The loop iterates until a fixed evaluation budget is exhausted, at which point the configuration with the highest validation accuracy is retained. That winning configuration is then rebuilt and retrained from scratch under the full training protocol, with early stopping on the validation set, and the resulting model is evaluated on the test set. This second phase serves a purpose beyond reporting a final number: it verifies whether the single-epoch fitness used during the search ranks configurations consistently with their fully trained performance, which is the assumption that makes the cheap proxy legitimate in the first place.

\subsection{Experimental Protocol}
\label{subsec:study1_protocol}

For each algorithm, the protocol is the same:

\begin{enumerate}
    \item \textbf{Independent repetitions.} Each algorithm is run independently $N_{\mathrm{val}} = 5$ times with an independently drawn seed per repetition. Each run executes the same fitness budget of $220$ evaluations, organized as $20$ individuals in the initial population plus $10$ generations of $20$ offspring for the population-based algorithms, and as $20$ random startup trials plus $200$ guided trials for \acrshort{bo}.

    \item \textbf{Control parameters.} The values of Table~\ref{tab:algorithm_params} are used and kept fixed across all repetitions. Algorithms that adapt their own control parameters (\acrshort{sade}, \acrshort{gwo}) are used in their library formulation with no additional tuning.

    \item \textbf{Final retraining.} The best configuration of each algorithm is retrained from scratch over $N_{\mathrm{re}} = 3$ repetitions with seeds $0$ to $2$, for up to $30$ epochs with early stopping (patience $5$) on the validation accuracy and restoration of the best weights.

    \item \textbf{Reporting.} For each algorithm the aggregated convergence behavior, the distribution of the final best fitness across the five repetitions, and a condensed summary table with the mean $\pm$ standard deviation as well as the median and range are reported.
\end{enumerate}

\subsection{Results}
\label{subsec:study1_results}

\paragraph{Convergence behavior.}
Figures~\ref{fig:study1_mnist_all_algorithms_convergence} and~\ref{fig:study1_cifar10_all_algorithms_convergence} report the best fitness (red) and the mean fitness (blue) of each generation, averaged over the $N_{\mathrm{val}}=5$ repetitions.

On MNIST the search operates against a ceiling: the best-of-generation curve is already near its final value at generation~$0$ and stays flat around $0.010$--$0.013$, since the initial random population alone already exceeds $99\%$ validation accuracy. The final best-fitness values are consequently tightly clustered (Table~\ref{tab:study1_mnist_comparison}, from $0.0082$ for \acrshort{ga} to $0.0108$ for \acrshort{gwo}), within roughly one standard deviation of each other. The discriminating signal is the population mean: \acrshort{ga}, \acrshort{sade} and \acrshort{gwo} drive it steadily downward, whereas \acrshort{aiwpso} and \acrshort{bo} keep a higher and noisier mean, consistent with more dispersed sampling, which is expected for \acrshort{bo} since it is not population-based.

On CIFAR-10 there is genuine room to improve and the profiles separate more clearly. From an initial best fitness around $0.50$, \acrshort{gwo} shows the steepest and most sustained descent, reaching the lowest and most consistent final fitness ($0.3919 \pm 0.0134$, Table~\ref{tab:study1_cifar10_comparison}), with both curves still decreasing at the last generation. \acrshort{sade} follows ($0.4013$), while \acrshort{aiwpso} improves early but stalls around generation~$4$ as its mean drifts back up, indicating premature convergence. \acrshort{ga} improves unevenly, with intermittent upticks between generations, and remains the weakest population-based method ($0.4276$), reversing its MNIST ranking. \acrshort{bo} reaches its best earliest, after $53.6$ evaluations on average, but finishes mid-pack ($0.4158$).

\begin{figure}[ht]
    \centering
    \includegraphics[width=\textwidth]{Chapters/img/mnist_1_all_algorithms_convergence_per_generation.pdf}
    \caption{MNIST: per-generation convergence of the five candidate algorithms, reported as mean $\pm$ one standard deviation over $N_{\mathrm{val}}=5$ independent runs per algorithm. Red: best fitness of each generation. Blue: mean fitness of each generation. Shaded bands represent $\pm 1$ standard deviation across the five runs. All panels share the same $Y$ axis.}
    \label{fig:study1_mnist_all_algorithms_convergence}
\end{figure}

\begin{figure}[ht]
    \centering
    \includegraphics[width=\textwidth]{Chapters/img/cifar_10_1_all_algorithms_convergence_per_generation.pdf}
    \caption{CIFAR-10: per-generation convergence of the five candidate algorithms, reported as mean $\pm$ one standard deviation over $N_{\mathrm{val}}=5$ independent runs per algorithm. Red: best fitness of each generation. Blue: mean fitness of each generation. Shaded bands represent $\pm 1$ standard deviation across the five runs. All panels share the same $Y$ axis.}
    \label{fig:study1_cifar10_all_algorithms_convergence}
\end{figure}

\paragraph{Distribution of the final solution.}
Figures~\ref{fig:study1_mnist_all_algorithms_boxplot} and~\ref{fig:study1_cifar10_all_algorithms_boxplot} show the distribution of the final best fitness across the five repetitions, which exposes the consistency of each algorithm beyond the aggregated means of Tables~\ref{tab:study1_mnist_comparison} and~\ref{tab:study1_cifar10_comparison}.

On MNIST (Figure~\ref{fig:study1_mnist_all_algorithms_boxplot}) all distributions are compressed into a narrow band, again reflecting the ceiling effect of the dataset. \acrshort{ga} attains the lowest mean ($0.0082$) but with a comparatively wide spread, while \acrshort{bo} is nearly as good ($0.0088$) and tighter. \acrshort{aiwpso} collapses almost entirely onto a single value ($0.0100 \pm 0.0001$), the most consistent behavior of all but at a mediocre fitness. \acrshort{gwo} is the worst on this dataset by both location and dispersion, with the highest median and the widest interquartile range.

On CIFAR-10 (Figure~\ref{fig:study1_cifar10_all_algorithms_boxplot}) the distributions separate clearly. \acrshort{gwo} yields the lowest mean ($0.3919$) and the lowest median while remaining among the tightest, confirming it as the most reliable optimizer on this search space. \acrshort{bo} is the most concentrated ($\pm 0.0131$) but centered on a higher value. The case of \acrshort{ga} is the most revealing: although it produced the single best individual run, its box sits high and it has both the worst median and the largest variance ($0.4276 \pm 0.0305$), so its competitive mean is an artifact of one fortunate run rather than typical behavior. \acrshort{sade} and \acrshort{aiwpso} occupy intermediate positions, with \acrshort{sade} showing a wider spread driven by its upper outlier.

Overall, the boxplots reinforce the convergence analysis and add a reliability dimension: \acrshort{gwo} and \acrshort{ga} swap places between the two datasets, but on the harder CIFAR-10 problem \acrshort{gwo} is both the best and one of the most stable methods, whereas \acrshort{ga} combines a strong best case with the poorest worst case and median.

\begin{figure}[ht]
    \centering
    \includegraphics[width=\textwidth]{Chapters/img/mnist_1_all_algorithms_boxplot_final_labeled.pdf}
    \caption{MNIST: distribution of the final best fitness reached by each candidate algorithm over $N_{\mathrm{val}}=5$ independent repetitions. The box gives the interquartile range, the red line the median, the white diamond the mean, the whiskers the minimum and maximum, and the blue points the individual runs (labelled $1$--$5$). All boxplots share the same $Y$ axis.}
    \label{fig:study1_mnist_all_algorithms_boxplot}
\end{figure}

\begin{figure}[ht]
    \centering
    \includegraphics[width=\textwidth]{Chapters/img/cifar_10_1_all_algorithms_boxplot_final_labeled.pdf}
    \caption{CIFAR-10: distribution of the final best fitness reached by each candidate algorithm over $N_{\mathrm{val}}=5$ independent repetitions. The box gives the interquartile range, the red line the median, the white diamond the mean, the whiskers the minimum and maximum, and the blue points the individual runs (labelled $1$--$5$). All boxplots share the same $Y$ axis.}
    \label{fig:study1_cifar10_all_algorithms_boxplot}
\end{figure}

\clearpage
\paragraph{Summary statistics.}
Tables~\ref{tab:study1_mnist_comparison} (MNIST) and~\ref{tab:study1_cifar10_comparison} (CIFAR-10) summarize the per-algorithm search statistics over the $N_{\mathrm{val}}=5$ independent runs of $220$ evaluations each. Each row reports the best fitness, the corresponding validation accuracy, the evaluation index at which the best was found, the wall-clock time, and the minimum, maximum and median across the five seeds.

Two observations follow from the wall-clock column. First, the cost of a run is dominated by the training of the candidates and not by the optimizer, which is why the totals are of the same order for every method. Second, the differences that do appear track the size of the architectures each optimizer favors rather than any intrinsic overhead, a point that Study~2 makes explicit by removing training from the loop entirely.

\begin{table}[h!]
    \centering
    \caption{MNIST: aggregated search results of the five candidate algorithms over $N_{\mathrm{val}} = 5$ independent repetitions of $220$ evaluations each. Values are reported as $\mathrm{mean} \pm \mathrm{std}$ unless otherwise stated.}
    \label{tab:study1_mnist_comparison}
    \resizebox{\textwidth}{!}{%
    \begin{tabular}{lcccccccc}
        \toprule
        \textbf{Algorithm} & \textbf{n\_runs} & \textbf{Best fit} & \textbf{Best val\_acc} & \textbf{Eval @ best} & \textbf{Time (s)} & \textbf{Min} & \textbf{Max} & \textbf{Median} \\
        \midrule
        GA      & 5 & $\mathbf{0.0082 \pm 0.0010}$ & 0.9918 $\pm$ 0.0010 & 178.2 $\pm$ 21.8 & 6863.6 $\pm$ 449.0 & 0.0067 & 0.0094 & 0.0083 \\
        SADE    & 5 & 0.0090 $\pm$ 0.0010 & 0.9910 $\pm$ 0.0010 & 166.6 $\pm$ 52.0 & 8647.8 $\pm$ 539.2 & 0.0076 & 0.0105 & 0.0089 \\
        AIW-PSO & 5 & 0.0100 $\pm$ 0.0001 & 0.9900 $\pm$ 0.0001 & 121.2 $\pm$ 60.5 & 6416.5 $\pm$ 205.8 & 0.0099 & 0.0101 & 0.0100 \\
        GWO     & 5 & 0.0108 $\pm$ 0.0009 & 0.9892 $\pm$ 0.0009 & 141.4 $\pm$ 71.2 & 9377.1 $\pm$ 484.3 & 0.0098 & 0.0118 & 0.0107 \\
        BO      & 5 & 0.0088 $\pm$ 0.0004 & 0.9912 $\pm$ 0.0004 & 98.2 $\pm$ 68.8  & 8112.5 $\pm$ 381.4 & 0.0083 & 0.0095 & 0.0087 \\
        \bottomrule
    \end{tabular}%
    }
\end{table}

\begin{table}[h!]
    \centering
    \caption{CIFAR-10: aggregated search results of the five candidate algorithms over $N_{\mathrm{val}} = 5$ independent repetitions of $220$ evaluations each. Values are reported as $\mathrm{mean} \pm \mathrm{std}$ unless otherwise stated.}
    \label{tab:study1_cifar10_comparison}
    \resizebox{\textwidth}{!}{%
    \begin{tabular}{lcccccccc}
        \toprule
        \textbf{Algorithm} & \textbf{n\_runs} & \textbf{Best fit} & \textbf{Best val\_acc} & \textbf{Eval @ best} & \textbf{Time (s)} & \textbf{Min} & \textbf{Max} & \textbf{Median} \\
        \midrule
        GA      & 5 & 0.4276 $\pm$ 0.0305 & 0.5724 $\pm$ 0.0305 & 190.0 $\pm$ 30.1 & 9015.6 $\pm$ 925.8  & 0.3696 & 0.4601 & 0.4356 \\
        SADE    & 5 & 0.4013 $\pm$ 0.0243 & 0.5987 $\pm$ 0.0243 & 154.6 $\pm$ 46.2 & 7362.2 $\pm$ 168.1  & 0.3789 & 0.4453 & 0.3905 \\
        AIW-PSO & 5 & 0.4126 $\pm$ 0.0170 & 0.5874 $\pm$ 0.0170 & 110.2 $\pm$ 31.4 & 7359.6 $\pm$ 724.3  & 0.3888 & 0.4325 & 0.4131 \\
        GWO     & 5 & $\mathbf{0.3919 \pm 0.0134}$ & 0.6081 $\pm$ 0.0134 & 177.8 $\pm$ 39.5 & 8061.1 $\pm$ 1191.4 & 0.3779 & 0.4097 & 0.3823 \\
        BO      & 5 & 0.4158 $\pm$ 0.0131 & 0.5842 $\pm$ 0.0131 & 53.6  $\pm$ 30.7 & 7097.2 $\pm$ 89.8   & 0.3963 & 0.4374 & 0.4159 \\
        \bottomrule
    \end{tabular}%
    }
\end{table}

\paragraph{Final retraining.}
For each algorithm, the best configuration found during the search was retrained from scratch under the full training protocol. Tables~\ref{tab:study1_mnist_retrain_summary} (MNIST) and~\ref{tab:study1_cifar10_retrain_summary} (CIFAR-10) report the resulting test accuracy, the epoch at which it was attained, the wall-clock training time and the number of trainable parameters of the chosen architecture. The corresponding learning trajectories are plotted in Figures~\ref{fig:study1_mnist_retrain_val_acc} and~\ref{fig:study1_cifar10_retrain_val_acc}, where each curve is the mean validation accuracy per epoch over $N_{\mathrm{re}}=3$ retraining seeds.

The retraining results support the legitimacy of the single-epoch proxy: on MNIST every selected configuration lands between $0.9939$ and $0.9954$ test accuracy, and on CIFAR-10 between $0.7886$ and $0.8207$, so the ordering of the configurations under full training is broadly consistent with the ordering under the cheap fitness. They also expose a dimension the search itself does not optimize: \acrshort{ga} and \acrshort{gwo} converge to the same compact model on each dataset, whereas \acrshort{aiwpso} and \acrshort{sade} select configurations more than twice that size for a comparable accuracy. This observation is what motivates making the cost an explicit objective in Study~2.

\begin{table}[h!]
    \centering
    \caption{MNIST: final retraining results over $N_{\mathrm{re}} = 3$ independent retraining seeds per configuration. \emph{HPO val\_acc} is the validation accuracy reported during the search phase (single-epoch training). \emph{Test acc} is the best accuracy reached during the full retraining. Values are $\mathrm{mean} \pm \mathrm{std}$ unless otherwise stated.}
    \label{tab:study1_mnist_retrain_summary}
    \resizebox{\textwidth}{!}{%
    \begin{tabular}{lccccc}
        \toprule
        \textbf{Algorithm} & \textbf{HPO val\_acc} & \textbf{Test acc} & \textbf{Best epoch} & \textbf{Train time (s)} & \textbf{\# Params} \\
        \midrule
        GA      & $\mathbf{0.9933}$ & $\mathbf{0.9954 \pm 0.0004}$ & 12.7 $\pm$ 5.0 & $\mathbf{352.3 \pm 99.4}$  & 437{,}098 \\
        SADE    & 0.9924 & 0.9949 $\pm$ 0.0005 & $\mathbf{10.7 \pm 2.9}$ & 383.2 $\pm$ 75.2  & 883{,}690 \\
        AIW-PSO & 0.9901 & 0.9939 $\pm$ 0.0028 & 11.0 $\pm$ 7.5 & 549.0 $\pm$ 258.0 & 1{,}094{,}378 \\
        GWO     & 0.9902 & 0.9951 $\pm$ 0.0001 & 12.7 $\pm$ 3.8 & 406.8 $\pm$ 88.2  & 437{,}098 \\
        BO      & 0.9917 & 0.9941 $\pm$ 0.0006 & 16.7 $\pm$ 6.7 & 521.1 $\pm$ 153.6 & 883{,}690 \\
        \bottomrule
    \end{tabular}%
    }
\end{table}

\begin{figure}[h!]
    \centering
    \includegraphics[width=\textwidth]{Chapters/img/mnist_1_retrain_val_accuracy_per_epoch.pdf}
    \caption{MNIST: validation accuracy per epoch for the best configuration of each algorithm, retrained from scratch with the full training protocol. Each curve is the mean over $N_{\mathrm{re}}=3$ independent retraining seeds, and the shaded band is $\pm 1$ standard deviation. Curves that end earlier than the others were stopped by the early-stopping criterion.}
    \label{fig:study1_mnist_retrain_val_acc}
\end{figure}

\begin{table}[h!]
    \centering
    \caption{CIFAR-10: final retraining results over $N_{\mathrm{re}} = 3$ independent retraining seeds per configuration. \emph{HPO val\_acc} is the validation accuracy reported during the search phase (single-epoch training). \emph{Val acc} and \emph{Test acc} are the best values reached during the full retraining. All variable values are $\mathrm{mean} \pm \mathrm{std}$.}
    \label{tab:study1_cifar10_retrain_summary}
    \resizebox{\textwidth}{!}{%
    \begin{tabular}{lcccccc}
        \toprule
        \textbf{Algorithm} & \textbf{HPO val\_acc} & \textbf{Val acc} & \textbf{Test acc} & \textbf{Best epoch} & \textbf{Train time (s)} & \textbf{\# Params} \\
        \midrule
        GA      & $\mathbf{0.6304}$ & 0.7942 $\pm$ 0.0058 & 0.7911 $\pm$ 0.0058 & 16.0 $\pm$ 5.6 & $\mathbf{379.7 \pm 98.1}$  & 552{,}362     \\
        SADE    & 0.6211 & 0.7938 $\pm$ 0.0098 & 0.7886 $\pm$ 0.0141 & 17.0 $\pm$ 8.5 & 495.3 $\pm$ 198.2 & 1{,}343{,}018 \\
        AIW-PSO & 0.6112 & 0.8116 $\pm$ 0.0058 & 0.8057 $\pm$ 0.0045 & 21.7 $\pm$ 7.2 & 825.5 $\pm$ 124.4 & 1{,}325{,}354 \\
        GWO     & 0.6221 & 0.8042 $\pm$ 0.0117 & 0.8019 $\pm$ 0.0064 & 19.3 $\pm$ 2.5 & 543.6 $\pm$ 55.9  & 552{,}362     \\
        BO      & 0.6037 & $\mathbf{0.8302 \pm 0.0017}$ & $\mathbf{0.8207 \pm 0.0053}$ & 23.0 $\pm$ 4.4 & 643.7 $\pm$ 79.4  & 1{,}343{,}018 \\
        \bottomrule
    \end{tabular}%
    }
\end{table}

\begin{figure}[h!]
    \centering
    \includegraphics[width=\textwidth]{Chapters/img/cifar_10_1_retrain_val_accuracy_per_epoch.pdf}
    \caption{CIFAR-10: validation accuracy per epoch for the best configuration of each algorithm, retrained from scratch with the full training protocol. Each curve is the mean over $N_{\mathrm{re}}=3$ independent retraining seeds, and the shaded band is $\pm 1$ standard deviation. Curves that end earlier than the others were stopped by the early-stopping criterion.}
    \label{fig:study1_cifar10_retrain_val_acc}
\end{figure}

\paragraph{Cross-dataset comparison.}
Figure~\ref{fig:study1_ranking_heatmap} renders the algorithm ranking compactly: the left panel ranks the five algorithms within each dataset (1 = best) and the right panel shows the mean final validation accuracy.

Two observations follow. First, the ranking is \emph{almost fully reversed} between the two datasets: \acrshort{ga} ranks first on MNIST and last on CIFAR-10, whereas \acrshort{gwo} does the opposite. Second, the saturation of MNIST compresses the five algorithms into a band of width $0.0026$ in validation accuracy, less than three standard deviations of the noisiest method, whereas CIFAR-10 spreads them across $0.0357$, an order of magnitude wider. The harder benchmark therefore acts as the actual discriminator between optimizers, and the absolute ranking on MNIST should be interpreted with the corresponding caution.

\begin{figure}[h!]
    \centering
    \includegraphics[width=\textwidth]{Chapters/img/hpo_ranking_heatmap.pdf}
    \caption{Cross-dataset algorithm comparison on the \acrshort{hpo} task. Left: ranking per dataset (1 = best). Right: mean final validation accuracy across the $N_{\mathrm{val}}=5$ seeds.}
    \label{fig:study1_ranking_heatmap}
\end{figure}

\subsection{Discussion}
\label{subsec:study1_discussion}

Three conclusions can be drawn from Study~1.

The difficulty of the benchmark determines how much it can discriminate between optimizers. MNIST saturates against a ceiling and compresses every algorithm into a near-indistinguishable band, in which secondary signals such as the population mean and the seed-to-seed dispersion carry more information than the headline best fitness. CIFAR-10 exposes clear and stable differences under the same budget and the same encoding.

\acrshort{gwo} is the strongest and most reliable optimizer on this categorical space, where the budget is small and each evaluation is expensive. \acrshort{sade} follows closely. Both required no configuration step, which is a practical advantage in exactly the setting where automated search is needed. \acrshort{ga} converges prematurely, and \acrshort{aiwpso} explores widely but exploits poorly, stalling after a few generations.

Reliability deserves to be treated as a first-class criterion rather than as a footnote to the mean. Under five repetitions the average can be misleading, as the \acrshort{ga} result on CIFAR-10 illustrates: it produced the single best individual run yet the worst median and the largest variance among the population-based methods, so its competitive mean is an artifact of one fortunate run. Reporting distributions, medians and ranges alongside means guards against this artifact, and this convention is kept throughout the remaining studies.

\clearpage
\section{Study 2 --- Neural Architecture Search on NAS-Bench-201}
\label{sec:study_2}

The second study replaces the categorical hyperparameter search of Section~\ref{sec:study_1} by a \acrfull{nas} problem. The \acrshort{cnn} template is no longer fixed: the optimizer chooses the operations inside a small repeating cell that defines the whole network.

To keep the cost tractable, the search is run on NAS-Bench-201~\cite{dong2020nasbench201}, a tabular benchmark that pre-trains and evaluates every architecture of a fixed cell-based search space. Each evaluation is therefore a database lookup rather than a training run, which has three consequences that shape the whole study. First, the evaluation budget can be raised well beyond what Study~1 allowed. Second, the objective is noise-free, so an algorithm is never penalized or rewarded by training variance. Third, and most importantly, both the global optimum and the exact Pareto front can be obtained by exhaustive enumeration of the $5^6 = 15{,}625$ architectures, and used as ground-truth references against which every method is measured.

The study is organized in three parts that share the same search space, the same pipeline and the same budget, and differ only in how the objective is formulated:

\begin{itemize}
    \item \textbf{Part~A} (Section~\ref{subsec:study2_single_objective}) treats \acrshort{nas} as a single-objective problem driven by accuracy alone;
    \item \textbf{Part~B} (Section~\ref{subsec:study2_scalarized}) makes the same optimizers cost-aware while keeping the problem single-objective, by scalarizing accuracy and computational cost into a weighted sum;
    \item \textbf{Part~C} (Section~\ref{subsec:study2_multiobjective}) optimizes accuracy and cost simultaneously with native Pareto optimizers, and compares the result against the scalarized baseline of Part~B.
\end{itemize}

\subsection{Datasets}
\label{subsec:study2_datasets}

NAS-Bench-201 ships with three image-classification datasets of increasing difficulty (Table~\ref{tab:study2_datasets}). All algorithms are run on each of them. CIFAR-10 also appears in Study~1, which allows a direct reading of how the same dataset behaves under two different search spaces.

\begin{table}[h]
    \centering
    \caption{Datasets bundled with NAS-Bench-201 and used in Study~2.}
    \label{tab:study2_datasets}
    \begin{tabular}{lccc}
        \toprule
        \textbf{Dataset} & \textbf{Labels} & \textbf{Input shape} & \textbf{Train / Test size} \\
        \midrule
        CIFAR-10       & 10  & $32\times32\times3$ & $50{,}000$ / $10{,}000$ \\
        CIFAR-100      & 100 & $32\times32\times3$ & $50{,}000$ / $10{,}000$ \\
        ImageNet16-120 & 120 & $16\times16\times3$ & $151{,}700$ / $6{,}000$ \\
        \bottomrule
    \end{tabular}
\end{table}

\subsection{Search Space and Encoding}
\label{subsec:study2_search_space}

Each candidate is a cell, modeled as a directed acyclic graph with four nodes and six edges. Every pair of nodes is connected exactly once in the forward direction, so node~$0$ is the cell's input, node~$3$ is its output, and nodes~$1$ and~$2$ are intermediate feature maps that accumulate the outputs of every earlier node. Each edge applies one of the five primitive operations listed in Table~\ref{tab:study2_search_space}.

A candidate solution is encoded as an integer vector of length six in $\{0, \dots, 4\}^{6}$, one entry per edge in the fixed order $0{\to}1$, $0{\to}2$, $0{\to}3$, $1{\to}2$, $1{\to}3$, $2{\to}3$, yielding $5^6 = 15{,}625$ unique cell topologies. The full network is built by stacking copies of the chosen cell separated by fixed reduction layers; this macro skeleton is not exposed to the optimizer.

\begin{table}[ht]
    \centering
    \caption{Search space of Study~2, with cardinality equal to $5^6 = 15{,}625$ unique cell topologies.}
    \label{tab:study2_search_space}
    \begin{tabular}{cll}
        \toprule
        \textbf{Index} & \textbf{Operation} & \textbf{Description} \\
        \midrule
        0 & \texttt{none}           & zeroize input \\
        1 & \texttt{skip\_connect}  & identity \\
        2 & \texttt{nor\_conv\_1x1} & $1\times1$ convolution $+$ BN $+$ ReLU \\
        3 & \texttt{nor\_conv\_3x3} & $3\times3$ convolution $+$ BN $+$ ReLU \\
        4 & \texttt{avg\_pool\_3x3} & $3\times3$ average pooling \\
        \bottomrule
    \end{tabular}
\end{table}

The six integer coordinates are kept native for the evolutionary operators. For the solvers that require a continuous domain, namely the differential-evolution and swarm updates and the Gaussian-process methods, the variables are relaxed to a continuous domain and rounded to the nearest operation index on evaluation, a standard practice for integer \acrshort{nas} spaces.

\subsection{Optimization Pipeline}
\label{subsec:study2_pipeline}

The pipeline of Section~\ref{sec:optimisation_pipeline} is reused with one change in the fitness step: the integer vector is decoded into the NAS-Bench-201 architecture string, the benchmark is queried, and the tabulated metrics of the corresponding architecture are returned. No network is trained at any point of the search. The exhaustive enumeration of the space, computed once per dataset, provides the reference optimum and the reference Pareto front. Figure~\ref{fig:nas_workflow} depicts the resulting workflow.

\begin{figure}[ht]
    \centering
    \includegraphics[width=\textwidth]{Chapters/img/nas_workflow.pdf}
    \caption{Workflow of the NAS-Bench-201 search pipeline. The search loop queries the tabular benchmark instead of training a \acrshort{cnn} at each iteration. The bottom section exploits the tractable size of the search space to compute the exhaustive optimum, used as a reference benchmark.}
    \label{fig:nas_workflow}
\end{figure}

Because no training happens inside the loop, the wall-clock time of a run measures the \emph{optimizer's own overhead} rather than the cost of evaluating candidates. This makes Study~2 the natural place to quantify how expensive each search strategy is in itself, a quantity that Study~1 could not isolate. Throughout this section the tabulated training cost of the queried architectures is therefore kept distinct from the search wall-clock, and only the latter is reported per algorithm.

\subsection{Problem Formulations and Quality Indicators}
\label{subsec:study2_formulations}

\paragraph{Accuracy-only formulation.}
In Part~A the goal is to maximize test accuracy, that is, to minimize

\begin{equation}
    \mathrm{fit}(\mathbf{x}) \;=\; 1 - \mathrm{acc}_{\mathrm{test}}(\mathbf{x}),
    \label{eq:fitness_study2}
\end{equation}

where $\mathrm{acc}_{\mathrm{test}}$ is the NAS-Bench-201 test accuracy of the architecture encoded by $\mathbf{x}$ after $200$ epochs of training.

\paragraph{Cost-aware scalarized formulation.}
Accuracy alone is rarely the deployment objective: a model must also be efficient in terms of size, latency or Floating-Point Operations (FLOPs). Part~B keeps the search single-objective while making it cost-aware, by scalarizing accuracy and FLOPs into a single weighted sum. The cost is the architecture FLOPs normalized to $[0,1]$ over the benchmark,

\begin{equation}
    \widehat{\mathrm{FLOPs}}(\mathbf{x}) = \frac{\mathrm{FLOPs}(\mathbf{x}) - \mathrm{FLOPs}_{\min}}{\mathrm{FLOPs}_{\max} - \mathrm{FLOPs}_{\min}},
    \label{eq:flops}
\end{equation}

so the two criteria to minimize are

\begin{equation}
    f_1(\mathbf{x}) = 1 - \mathrm{acc}_{\mathrm{test}}(\mathbf{x}), \qquad f_2(\mathbf{x}) = \widehat{\mathrm{FLOPs}}(\mathbf{x}),
    \label{eq:objectives}
\end{equation}

and, for a fixed weight $\alpha$, the optimizer minimizes

\begin{equation}
    g_\alpha(\mathbf{x}) = \alpha\, f_1(\mathbf{x}) + (1-\alpha)\, f_2(\mathbf{x}).
    \label{eq:scalar}
\end{equation}

A larger $\alpha$ favors accuracy and a smaller one favors cost. Sweeping $\alpha$ over $\{0.1, \dots, 0.9\}$ yields nine independent runs per optimizer, whose non-dominated best architectures approximate the accuracy--cost trade-off and serve as the single-objective baseline for Part~C.

\paragraph{Multi-objective formulation.}
Part~C optimizes the two criteria of Equation~\eqref{eq:objectives} simultaneously rather than collapsing them into a weighted sum. A solution $\mathbf{x}$ dominates $\mathbf{y}$, written $\mathbf{x} \prec \mathbf{y}$, when it is no worse in both objectives and strictly better in at least one,

\begin{equation}
    \mathbf{x} \prec \mathbf{y} \iff f_i(\mathbf{x}) \le f_i(\mathbf{y})\ \ \forall i \in \{1,2\} \ \ \wedge\ \ \exists\, j \in \{1,2\}: f_j(\mathbf{x}) < f_j(\mathbf{y}),
    \label{eq:dominance}
\end{equation}

and the goal is to approximate the Pareto front, the set of non-dominated architectures.

The quality of an approximation $S$ is measured by two complementary indicators. The \emph{hypervolume} is the area of objective space dominated by $S$ and bounded by a reference point $R = (1.1, 1.1)$ placed beyond the worst attainable values~\cite{zitzler1999hv},

\begin{equation}
    \mathrm{HV}(S) = \Lambda\!\left( \bigcup_{\mathbf{x}\in S} \big[f_1(\mathbf{x}), R_1\big] \times \big[f_2(\mathbf{x}), R_2\big] \right),
    \label{eq:hv}
\end{equation}

where $\Lambda$ denotes the Lebesgue measure. Since the benchmark is tabular, the exact Pareto front of each dataset is obtained by enumeration and yields the maximum attainable hypervolume $\mathrm{HV}^{\star}$, so the gap $\mathrm{HV}^{\star} - \mathrm{HV}$ measures how far an algorithm is from the optimum. The hypervolume jointly rewards convergence and diversity, but it can be dominated by a dense cluster of points in a cheap region of the front; for this reason the Inverted Generational Distance Plus (IGD$^{+}$)~\cite{ishibuchi2015igdplus} is also reported, since it explicitly penalizes incomplete coverage of the true front.

\subsection{Experimental Protocol}
\label{subsec:study2_protocol}

The protocol is identical across algorithms within each formulation:

\begin{enumerate}
    \item \textbf{Accuracy-only (Part~A).} Each (algorithm, dataset) pair is run $N_{\mathrm{val}} = 5$ times with seeds $0$ to $4$, matched across algorithms, under a budget of $1020$ queries, organized as $20$ initial individuals plus $50$ generations of $20$ offspring, or, for \acrshort{bo}, as $20$ random startup trials plus $1000$ guided trials.

    \item \textbf{Scalarized (Part~B).} The same budget is spent once per weight $\alpha \in \{0.1, \dots, 0.9\}$, and each configuration is repeated over $N_{\mathrm{val}} = 3$ seeds. The assembled scalarized front therefore consumes nine times the per-run budget of a native multi-objective search, $9 \times 1020$ against $1020$ evaluations, so the comparison in Part~C is deliberately generous to scalarization.

    \item \textbf{Multi-objective (Part~C).} Each native optimizer is run with $N_{\mathrm{val}} = 3$ independent seeds under the same $1020$-evaluation budget.

    \item \textbf{References and reporting.} The upper bound of each dataset and its exact Pareto front are computed by exhaustive enumeration of all $15{,}625$ architectures. For every setting the convergence behavior, the distribution of the final result across repetitions, and condensed summary statistics are reported.
\end{enumerate}

\subsection{Part A --- Single-Objective NAS}
\label{subsec:study2_single_objective}

\paragraph{Convergence behavior.}
Figure~\ref{fig:study2_conv_all} reports the running best test accuracy per generation on the three datasets, with one panel per dataset.

On CIFAR-10 and CIFAR-100 every algorithm except \acrshort{aiwpso} reaches the upper bound. \acrshort{ga} is the fastest to climb but plateaus marginally below the optimum, an early symptom of premature convergence. \acrshort{sade}, \acrshort{gwo} and \acrshort{bo} reach the upper bound by generation $\sim\!25$. \acrshort{aiwpso} stalls below all the others.

On ImageNet16-120 the curves separate clearly. \acrshort{bo} converges to the upper bound by generation $\sim\!20$. \acrshort{gwo} and \acrshort{sade} close most of the gap but stabilize $\approx 10^{-3}$ below it. \acrshort{ga} and \acrshort{aiwpso} converge to a clearly inferior plateau.

\begin{figure}[ht]
    \centering
    \includegraphics[width=\textwidth]{Chapters/img/nas_convergence_per_generation_all.pdf}
    \caption{NAS-Bench-201 convergence on the three datasets. Running best test accuracy per generation, mean $\pm$ one standard deviation over five seeds. The dashed line is the exhaustive upper bound of each dataset.}
    \label{fig:study2_conv_all}
\end{figure}

\paragraph{Distance to the upper bound.}
Because the optimum is known by exhaustive enumeration, the gap $\mathrm{ub} - \mathrm{best}_{\mathrm{sofar}}$ can be plotted on a logarithmic scale (Figure~\ref{fig:study2_gap_all}), which exposes differences that the linear convergence plots hide.

Three regimes appear consistently across datasets. \acrshort{bo} drives the gap to the floor of $10^{-6}$, which is machine-precision zero. \acrshort{sade} and \acrshort{gwo} descend by two further orders of magnitude after generation $\sim\!15$. \acrshort{ga} plateaus around $10^{-4}$ on CIFAR-10 and CIFAR-100 and around $5 \times 10^{-3}$ on ImageNet16-120. \acrshort{aiwpso} never descends below $\sim\!2 \times 10^{-3}$.

\begin{figure}[ht]
    \centering
    \includegraphics[width=\textwidth]{Chapters/img/nas_gap_log_all.pdf}
    \caption{NAS-Bench-201 distance to the exhaustive upper bound on a logarithmic scale, on the three datasets (mean over five seeds). The floor at $10^{-6}$ signals convergence to the optimum; lower is better.}
    \label{fig:study2_gap_all}
\end{figure}

\paragraph{Distribution of the final solution.}
Figure~\ref{fig:study2_box_all} shows the distribution of the final best test accuracy across the five seeds, with one panel per dataset.

On CIFAR-10 and CIFAR-100, \acrshort{ga}, \acrshort{sade}, \acrshort{gwo} and \acrshort{bo} collapse onto the upper bound, and \acrshort{aiwpso} is the only algorithm whose box is wide and offset, showing strong seed-to-seed instability. On ImageNet16-120 all five distributions separate: \acrshort{bo} is the most reliable, with all five runs at the upper bound; \acrshort{gwo} and \acrshort{sade} have a few runs at the optimum and the rest within $10^{-3}$ of it; \acrshort{ga} and \acrshort{aiwpso} show the largest spreads and the lowest medians.

\begin{figure}[ht]
    \centering
    \includegraphics[width=\textwidth]{Chapters/img/nas_boxplot_final_all.pdf}
    \caption{NAS-Bench-201 distribution of the final best test accuracy across five seeds per algorithm, on the three datasets. The dashed line is the exhaustive upper bound of each dataset.}
    \label{fig:study2_box_all}
\end{figure}

\paragraph{Summary statistics and cost of the selected architectures.}
Table~\ref{tab:study2_summary} consolidates the final accuracies and the gap to the exhaustive upper bound across the three datasets. Table~\ref{tab:study2_arch_cost} reports the cost of the architecture each algorithm selected, in FLOPs and in number of parameters. On CIFAR-10 and CIFAR-100 the methods that reach the upper bound all converge to the same architecture, which is why their cost figures coincide and their standard deviations vanish; \acrshort{aiwpso} settles on less computationally demanding but less accurate cells.

\begin{table}[ht]
    \centering
    \caption{NAS-Bench-201 accuracy-only results: mean final test accuracy $\pm$ std over $N_{\mathrm{val}}=5$ seeds, with the gap to the exhaustive upper bound in parentheses. The last row reports the upper bound itself for reference.}
    \label{tab:study2_summary}
    \resizebox{\textwidth}{!}{%
    \begin{tabular}{lccc}
        \toprule
        \textbf{Algorithm} & \textbf{CIFAR-10} & \textbf{CIFAR-100} & \textbf{ImageNet16-120} \\
        \midrule
        GA          & $0.9436 \pm 0.0003$ \, ($1\!\times\!10^{-4}$)   & $\mathbf{0.7351 \pm 0.0000}$ \, ($0$)          & $0.4686 \pm 0.0041$ \, ($4.5\!\times\!10^{-3}$) \\
        SADE        & $\mathbf{0.9437 \pm 0.0000}$ \, ($0$)           & $\mathbf{0.7351 \pm 0.0000}$ \, ($0$)          & $0.4719 \pm 0.0027$ \, ($1.2\!\times\!10^{-3}$) \\
        AIW-PSO     & $0.9417 \pm 0.0028$ \, ($2.0\!\times\!10^{-3}$) & $0.7327 \pm 0.0036$ \, ($2.5\!\times\!10^{-3}$) & $0.4674 \pm 0.0047$ \, ($5.7\!\times\!10^{-3}$) \\
        GWO         & $\mathbf{0.9437 \pm 0.0000}$ \, ($0$)           & $\mathbf{0.7351 \pm 0.0000}$ \, ($0$)          & $0.4722 \pm 0.0021$ \, ($9.0\!\times\!10^{-4}$) \\
        BO          & $\mathbf{0.9437 \pm 0.0000}$ \, ($0$)           & $\mathbf{0.7351 \pm 0.0000}$ \, ($0$)          & $\mathbf{0.4731 \pm 0.0000}$ \, ($0$)           \\
        \midrule
        Upper bound & $0.9437$ & $0.7351$ & $0.4731$ \\
        \bottomrule
    \end{tabular}%
    }
\end{table}

\begin{table}[ht]
    \centering
    \caption{Cost of the best architecture selected by each algorithm in the accuracy-only formulation, as \#FLOPs (M) and \#Params (M), $\pm$ std over $N_{\mathrm{val}}=5$ seeds. A zero standard deviation, omitted here, indicates that every seed converged to the same architecture.}
    \label{tab:study2_arch_cost}
    \resizebox{\textwidth}{!}{%
    \begin{tabular}{lcccccc}
        \toprule
        & \multicolumn{2}{c}{\textbf{CIFAR-10}} & \multicolumn{2}{c}{\textbf{CIFAR-100}} & \multicolumn{2}{c}{\textbf{ImageNet16-120}} \\
        \cmidrule(lr){2-3} \cmidrule(lr){4-5} \cmidrule(lr){6-7}
        \textbf{Algorithm} & \textbf{\#FLOPs} & \textbf{\#Params} & \textbf{\#FLOPs} & \textbf{\#Params} & \textbf{\#FLOPs} & \textbf{\#Params} \\
        \midrule
        GA      & $140.7 \pm 28.1$ & $0.99 \pm 0.19$ & $184.7$ & $1.29$ & $27.7 \pm 4.3$ & $0.79 \pm 0.12$ \\
        SADE    & $153.3$ & $1.07$ & $184.7$ & $1.29$ & $32.0 \pm 3.5$ & $0.91 \pm 0.10$ \\
        AIW-PSO & $165.9 \pm 30.4$ & $1.16 \pm 0.21$ & $158.8 \pm 35.6$ & $1.12 \pm 0.24$ & $33.2 \pm 4.7$ & $0.94 \pm 0.13$ \\
        GWO     & $153.3$ & $1.07$ & $184.7$ & $1.29$ & $31.8 \pm 3.1$ & $0.90 \pm 0.08$ \\
        BO      & $153.3$ & $1.07$ & $184.7$ & $1.29$ & $30.5$ & $0.87$ \\
        \bottomrule
    \end{tabular}%
    }
\end{table}

\paragraph{Architectural diversity.}
Figure~\ref{fig:study2_n_unique} reports the mean number of distinct architectures visited per run, out of the $15{,}625$ possible. The values are governed mainly by the algorithm rather than by the dataset: \acrshort{ga} stays around $120$ unique cells, \acrshort{sade} around $300$, \acrshort{gwo} around $400$, \acrshort{aiwpso} around $500$, and \acrshort{bo} reaches the full $1020$ queries without a single repetition. Across the three datasets each algorithm stays within the same band, with only a mild tendency for \acrshort{sade} and \acrshort{aiwpso} to explore more on ImageNet16-120, in each case within one standard deviation.

The figure is most informative when read jointly with the gap analysis, because it exposes the mechanism behind it: \acrshort{bo} converges to the optimum precisely because it never re-evaluates the same point, \acrshort{ga}'s narrow diversity coincides with its premature convergence, and \acrshort{aiwpso}'s high diversity is undirected and does not translate into accuracy.

\begin{figure}[ht]
    \centering
    \includegraphics[width=0.8\textwidth]{Chapters/img/nas_n_unique.pdf}
    \caption{Architectural diversity: mean number of unique cells queried per run for each (algorithm, dataset) pair, out of the $15{,}625$ architectures in NAS-Bench-201. Error bars are $\pm 1$ standard deviation over five seeds.}
    \label{fig:study2_n_unique}
\end{figure}

\paragraph{Search cost.}
Since the benchmark is tabular, the wall-clock of a run reflects optimizer overhead and not network training. Table~\ref{tab:study2_time} reports it per algorithm, and the differences span orders of magnitude: the four evolutionary and swarm methods complete a $1020$-evaluation run in about $6$ seconds, whereas \acrshort{bo} requires roughly $2700$ seconds, about $450\times$ more, because its Gaussian-process surrogate must be refit at every query and scales cubically with the number of evaluations. The tabulated training cost of the queried architectures is by contrast comparable across methods. \acrshort{bo} therefore reaches the optimum, but at a search cost two orders of magnitude above the evolutionary alternatives.

\begin{table}[ht]
    \centering
    \caption{Mean search wall-clock per run, in seconds, for the accuracy-only optimizers of Part~A ($N_{\mathrm{val}}=5$ seeds) and the native multi-objective optimizers of Part~C ($N_{\mathrm{val}}=3$ seeds). Since the benchmark is tabular, this reflects optimizer overhead and not network training; the Gaussian-process methods (\acrshort{bo} and qEHVI) scale cubically with the number of evaluations.}
    \label{tab:study2_time}
    \begin{tabular}{lccc}
        \toprule
        \textbf{Algorithm} & \textbf{CIFAR-10} & \textbf{CIFAR-100} & \textbf{ImageNet16-120} \\
        \midrule
        \multicolumn{4}{l}{\textit{Part A --- single-objective}} \\
        GA      & $5.9$    & $6.3$    & $6.0$    \\
        SADE    & $5.6$    & $6.4$    & $6.1$    \\
        AIW-PSO & $8.4$    & $6.4$    & $6.9$    \\
        GWO     & $7.5$    & $6.2$    & $6.2$    \\
        BO      & $2716.3$ & $2747.7$ & $2697.1$ \\
        \midrule
        \multicolumn{4}{l}{\textit{Part C --- multi-objective, native}} \\
        NSGA-II & $7.0$    & $6.9$    & $6.6$    \\
        GDE3    & $12.2$   & $9.5$    & $9.0$    \\
        qEHVI   & $3234.6$ & $2322.4$ & $2564.9$ \\
        \bottomrule
    \end{tabular}
\end{table}

\paragraph{Cross-dataset summary.}
Figure~\ref{fig:study2_heatmap} renders the per-dataset ranking as a heatmap. The relative ordering of algorithms is consistent only at the extremes: \acrshort{bo} is first or tied on every dataset and \acrshort{aiwpso} is last on every dataset. The middle of the ranking moves with the problem: \acrshort{ga} is first on CIFAR-100, second on CIFAR-10, but only fourth on ImageNet16-120, while \acrshort{sade} and \acrshort{gwo} are tied for second on CIFAR-100 and separate clearly on ImageNet16-120.

\begin{figure}[ht]
    \centering
    \includegraphics[width=\textwidth]{Chapters/img/nas_ranking_heatmap.pdf}
    \caption{Cross-dataset algorithm comparison on \acrshort{nas}. Left: ranking per dataset (1 = best). Right: mean final test accuracy.}
    \label{fig:study2_heatmap}
\end{figure}

CIFAR-10 and CIFAR-100 are too easy to discriminate the algorithms, since every method except \acrshort{aiwpso} reaches the optimum; ImageNet16-120 separates them. This mirrors, on a different search space, the ceiling effect that MNIST produced in Study~1.

\clearpage
\subsection{Part B --- Cost-Aware Scalarized NAS}
\label{subsec:study2_scalarized}

Part~A optimizes accuracy alone, and Table~\ref{tab:study2_arch_cost} already shows that the architectures it selects are among the most expensive in the benchmark. Part~B applies the same five optimizers to the scalarized objective of Equation~\eqref{eq:scalar}, which keeps the search single-objective while making it cost-aware, and sweeps the weight $\alpha$ from $0.1$ (cost-oriented) to $0.9$ (accuracy-oriented).

\paragraph{Assembled fronts.}
Figure~\ref{fig:study2_scalar_fronts} shows, for each optimizer, the non-dominated set assembled from its nine $\alpha$ runs, overlaid on the exact front. The pattern is the same for all five methods: each concentrates in the low-cost region and stops well short of the high-accuracy, high-FLOPs end of the front. Two causes contribute. One is intrinsic to a weighted sum, which cannot recover non-convex portions of a front. The other is the cap at $\alpha = 0.9$, which leaves a residual tenth of the weight always rewarding cost, so no run is ever driven purely by accuracy.

\begin{figure}[!ht]
    \centering
    \includegraphics[width=\textwidth]{Chapters/img/nas_scalar_fronts_a.pdf}\\[2mm]
    \includegraphics[width=\textwidth]{Chapters/img/nas_scalar_fronts_b.pdf}
    \caption{Scalarized formulation: non-dominated set assembled by each optimizer from its nine $\alpha$ runs, shown per optimizer and overlaid on the exact front (dashed). Each panel contains the three datasets, which occupy distinct accuracy bands. All optimizers concentrate in the low-cost region and stop short of the high-accuracy end of the front.}
    \label{fig:study2_scalar_fronts}
\end{figure}

\paragraph{Effect of the weight.}
Figure~\ref{fig:study2_alpha_sweep} plots the accuracy--FLOPs trajectory as $\alpha$ is swept, and Table~\ref{tab:study2_alpha_sweep} documents the exact architecture returned at every $\alpha$. The five trajectories almost coincide, which indicates that the scalarized optimizers explore essentially the same trade-off and differ mainly in how far they push along it. In other words, once the objective is scalarized, the choice of optimizer matters far less than the choice of $\alpha$.

\begin{figure}[!ht]
    \centering
    \includegraphics[width=\textwidth]{Chapters/img/nas_alpha_sweep.pdf}
    \caption{Scalarization $\alpha$-sweep: accuracy versus \#FLOPs of the architecture found by each scalarized algorithm as the weight $\alpha$ is swept from $0.1$ (cost-oriented) to $0.9$ (accuracy-oriented), per dataset. The five trajectories almost coincide, showing that the scalarized optimizers trace essentially the same accuracy--cost front.}
    \label{fig:study2_alpha_sweep}
\end{figure}

\begin{table}[!ht]
    \centering
    \caption{Scalarization $\alpha$-sweep: architecture returned by each scalarized algorithm for every weight $\alpha$, reported as Top-1 accuracy\,/\,\#FLOPs~(M)\,/\,\#Params~(M), mean over three seeds. Larger $\alpha$ favors accuracy, smaller $\alpha$ favors cost.}
    \label{tab:study2_alpha_sweep}
    \setlength{\tabcolsep}{3pt}
    \resizebox{\textwidth}{!}{%
    \begin{tabular}{clccccc}
        \toprule
        \textbf{Dataset} & $\mathbf{\alpha}$ & \textbf{GA} & \textbf{SADE} & \textbf{GWO} & \textbf{AIW-PSO} & \textbf{BO} \\
        \midrule
        \multirow{9}{*}{CIFAR-10} & 0.1 & 0.868\,/\,7.8\,/\,0.07 & 0.868\,/\,7.8\,/\,0.07 & 0.868\,/\,7.8\,/\,0.07 & 0.866\,/\,7.8\,/\,0.07 & 0.867\,/\,7.8\,/\,0.07 \\
         & 0.2 & 0.868\,/\,7.8\,/\,0.07 & 0.868\,/\,7.8\,/\,0.07 & 0.867\,/\,7.8\,/\,0.07 & 0.866\,/\,7.8\,/\,0.07 & 0.868\,/\,7.8\,/\,0.07 \\
         & 0.3 & 0.868\,/\,7.8\,/\,0.07 & 0.868\,/\,7.8\,/\,0.07 & 0.867\,/\,7.8\,/\,0.07 & 0.876\,/\,9.1\,/\,0.08 & 0.868\,/\,7.8\,/\,0.07 \\
         & 0.4 & 0.897\,/\,11.7\,/\,0.10 & 0.899\,/\,11.7\,/\,0.10 & 0.897\,/\,11.7\,/\,0.10 & 0.896\,/\,11.7\,/\,0.10 & 0.899\,/\,11.7\,/\,0.10 \\
         & 0.5 & 0.899\,/\,11.7\,/\,0.10 & 0.899\,/\,11.7\,/\,0.10 & 0.898\,/\,11.7\,/\,0.10 & 0.897\,/\,11.7\,/\,0.10 & 0.899\,/\,11.7\,/\,0.10 \\
         & 0.6 & 0.898\,/\,11.7\,/\,0.10 & 0.903\,/\,13.0\,/\,0.11 & 0.899\,/\,11.7\,/\,0.10 & 0.898\,/\,11.7\,/\,0.10 & 0.899\,/\,11.7\,/\,0.10 \\
         & 0.7 & 0.908\,/\,15.6\,/\,0.13 & 0.908\,/\,15.6\,/\,0.13 & 0.904\,/\,14.3\,/\,0.12 & 0.906\,/\,15.6\,/\,0.13 & 0.905\,/\,14.3\,/\,0.12 \\
         & 0.8 & 0.910\,/\,17.0\,/\,0.14 & 0.911\,/\,17.0\,/\,0.14 & 0.910\,/\,17.0\,/\,0.14 & 0.910\,/\,15.6\,/\,0.13 & 0.909\,/\,15.6\,/\,0.13 \\
         & 0.9 & 0.932\,/\,48.4\,/\,0.35 & 0.935\,/\,49.7\,/\,0.36 & 0.935\,/\,49.7\,/\,0.36 & 0.934\,/\,49.7\,/\,0.36 & 0.935\,/\,49.7\,/\,0.36 \\
        \midrule
        \multirow{9}{*}{CIFAR-100} & 0.1 & 0.589\,/\,7.8\,/\,0.08 & 0.587\,/\,7.8\,/\,0.08 & 0.588\,/\,7.8\,/\,0.08 & 0.560\,/\,7.8\,/\,0.08 & 0.588\,/\,7.8\,/\,0.08 \\
         & 0.2 & 0.586\,/\,7.8\,/\,0.08 & 0.588\,/\,7.8\,/\,0.08 & 0.585\,/\,7.8\,/\,0.08 & 0.602\,/\,9.1\,/\,0.09 & 0.588\,/\,7.8\,/\,0.08 \\
         & 0.3 & 0.650\,/\,11.7\,/\,0.11 & 0.646\,/\,11.7\,/\,0.11 & 0.644\,/\,11.7\,/\,0.11 & 0.654\,/\,13.0\,/\,0.12 & 0.645\,/\,11.7\,/\,0.11 \\
         & 0.4 & 0.647\,/\,11.7\,/\,0.11 & 0.647\,/\,11.7\,/\,0.11 & 0.647\,/\,11.7\,/\,0.11 & 0.645\,/\,11.7\,/\,0.11 & 0.646\,/\,11.7\,/\,0.11 \\
         & 0.5 & 0.669\,/\,15.7\,/\,0.14 & 0.670\,/\,15.7\,/\,0.14 & 0.670\,/\,15.7\,/\,0.14 & 0.659\,/\,14.3\,/\,0.13 & 0.671\,/\,15.7\,/\,0.14 \\
         & 0.6 & 0.670\,/\,15.7\,/\,0.14 & 0.669\,/\,15.7\,/\,0.14 & 0.671\,/\,15.7\,/\,0.14 & 0.671\,/\,17.0\,/\,0.14 & 0.670\,/\,15.7\,/\,0.14 \\
         & 0.7 & 0.679\,/\,20.9\,/\,0.17 & 0.675\,/\,18.3\,/\,0.15 & 0.674\,/\,18.3\,/\,0.15 & 0.673\,/\,18.3\,/\,0.15 & 0.674\,/\,18.3\,/\,0.15 \\
         & 0.8 & 0.683\,/\,23.5\,/\,0.19 & 0.684\,/\,23.5\,/\,0.19 & 0.684\,/\,23.5\,/\,0.19 & 0.684\,/\,23.5\,/\,0.19 & 0.682\,/\,22.2\,/\,0.18 \\
         & 0.9 & 0.716\,/\,55.0\,/\,0.41 & 0.714\,/\,53.7\,/\,0.40 & 0.716\,/\,55.0\,/\,0.41 & 0.709\,/\,56.3\,/\,0.42 & 0.715\,/\,56.3\,/\,0.42 \\
        \midrule
        \multirow{9}{*}{ImageNet16-120} & 0.1 & 0.286\,/\,2.0\,/\,0.08 & 0.286\,/\,2.0\,/\,0.08 & 0.285\,/\,2.0\,/\,0.08 & 0.283\,/\,2.0\,/\,0.08 & 0.285\,/\,2.0\,/\,0.08 \\
         & 0.2 & 0.284\,/\,2.0\,/\,0.08 & 0.286\,/\,2.0\,/\,0.08 & 0.285\,/\,2.0\,/\,0.08 & 0.284\,/\,2.0\,/\,0.08 & 0.286\,/\,2.0\,/\,0.08 \\
         & 0.3 & 0.341\,/\,2.9\,/\,0.11 & 0.341\,/\,2.9\,/\,0.11 & 0.340\,/\,2.9\,/\,0.11 & 0.341\,/\,2.9\,/\,0.11 & 0.357\,/\,3.3\,/\,0.12 \\
         & 0.4 & 0.368\,/\,3.6\,/\,0.13 & 0.383\,/\,3.9\,/\,0.14 & 0.383\,/\,3.9\,/\,0.14 & 0.366\,/\,3.9\,/\,0.14 & 0.383\,/\,3.9\,/\,0.14 \\
         & 0.5 & 0.398\,/\,4.9\,/\,0.16 & 0.384\,/\,3.9\,/\,0.14 & 0.382\,/\,3.9\,/\,0.14 & 0.379\,/\,3.9\,/\,0.14 & 0.384\,/\,3.9\,/\,0.14 \\
         & 0.6 & 0.398\,/\,4.9\,/\,0.16 & 0.402\,/\,5.2\,/\,0.17 & 0.394\,/\,4.6\,/\,0.16 & 0.389\,/\,4.9\,/\,0.16 & 0.401\,/\,5.2\,/\,0.17 \\
         & 0.7 & 0.407\,/\,5.6\,/\,0.18 & 0.406\,/\,5.9\,/\,0.19 & 0.398\,/\,5.2\,/\,0.17 & 0.400\,/\,5.9\,/\,0.19 & 0.407\,/\,5.6\,/\,0.18 \\
         & 0.8 & 0.454\,/\,14.1\,/\,0.42 & 0.458\,/\,14.4\,/\,0.43 & 0.449\,/\,13.7\,/\,0.41 & 0.454\,/\,14.4\,/\,0.43 & 0.454\,/\,13.7\,/\,0.41 \\
         & 0.9 & 0.454\,/\,14.4\,/\,0.43 & 0.460\,/\,14.7\,/\,0.44 & 0.460\,/\,14.7\,/\,0.44 & 0.453\,/\,14.7\,/\,0.44 & 0.458\,/\,14.4\,/\,0.43 \\
        \bottomrule
    \end{tabular}}
\end{table}

\paragraph{Quality of the assembled fronts.}
Table~\ref{tab:study2_scalar_hv} summarizes the quality of the assembled fronts as their hypervolume and gap to $\mathrm{HV}^{\star}$. The values are close to one another, confirming the reading of Figure~\ref{fig:study2_alpha_sweep}, and all of them fall short of what native multi-objective search achieves in Part~C, despite consuming nine times the budget.

\begin{table}[!ht]
    \centering
    \caption{Scalarized formulation: mean hypervolume $\pm$ std over $N_{\mathrm{val}}=3$ seeds of the front assembled from the nine $\alpha$ runs, with the gap to the true Pareto front $(\mathrm{HV}^{\star}-\mathrm{HV})$ in parentheses. The reference point is $R=(1.1,1.1)$. The true-front hypervolume $\mathrm{HV}^{\star}$ is given in the last row. The best value per dataset is in bold.}
    \label{tab:study2_scalar_hv}
    \resizebox{\textwidth}{!}{%
    \begin{tabular}{lccc}
        \toprule
        \textbf{Algorithm} & \textbf{CIFAR-10} & \textbf{CIFAR-100} & \textbf{ImageNet16-120} \\
        \midrule
        GA      & $1.1296 \pm 0.0048$ \, ($1.1\!\times\!10^{-2}$) & $0.8871 \pm 0.0004$ \, ($1.4\!\times\!10^{-2}$) & $0.5960 \pm 0.0023$ \, ($1.5\!\times\!10^{-2}$) \\
        SADE    & $\mathbf{1.1328 \pm 0.0009}$ \, ($7.3\!\times\!10^{-3}$) & $0.8858 \pm 0.0028$ \, ($1.5\!\times\!10^{-2}$) & $\mathbf{0.5989 \pm 0.0006}$ \, ($1.2\!\times\!10^{-2}$) \\
        GWO     & $1.1325 \pm 0.0009$ \, ($7.6\!\times\!10^{-3}$) & $\mathbf{0.8879 \pm 0.0000}$ \, ($1.3\!\times\!10^{-2}$) & $\mathbf{0.5989 \pm 0.0006}$ \, ($1.2\!\times\!10^{-2}$) \\
        AIW-PSO & $1.1315 \pm 0.0009$ \, ($8.6\!\times\!10^{-3}$) & $0.8809 \pm 0.0013$ \, ($2.0\!\times\!10^{-2}$) & $0.5931 \pm 0.0002$ \, ($1.8\!\times\!10^{-2}$) \\
        BO      & $1.1326 \pm 0.0007$ \, ($7.5\!\times\!10^{-3}$) & $0.8861 \pm 0.0020$ \, ($1.5\!\times\!10^{-2}$) & $0.5978 \pm 0.0026$ \, ($1.3\!\times\!10^{-2}$) \\
        \midrule
        True Pareto front $\mathrm{HV}^{\star}$ & $1.1401$ & $0.9008$ & $0.6107$ \\
        \bottomrule
    \end{tabular}}
\end{table}

\clearpage
\subsection{Part C --- Native Multi-Objective NAS}
\label{subsec:study2_multiobjective}

Part~C optimizes accuracy and cost simultaneously with the three native Pareto optimizers introduced in Section~\ref{subsec:algorithms}, under the same $1020$-evaluation budget, and compares them against the scalarized baseline of Part~B.

\paragraph{Hypervolume.}
Table~\ref{tab:study2_mo_hv} reports the hypervolume of the three native optimizers against the exact Pareto front. GDE3 attains the highest hypervolume and the smallest gap to the true front on all three datasets ($4.1 \times 10^{-4}$, $1.5 \times 10^{-3}$ and $1.3 \times 10^{-3}$), recovering both extremes of the front, including the most accurate architecture of each dataset ($0.9437$, $0.7351$, $0.4731$) and the least computationally demanding one. A Wilcoxon signed-rank test~\cite{demsar2006statistical} over the nine (dataset, seed) configurations confirms GDE3 as significantly the best, with a higher hypervolume than every other method in all nine cells ($p \approx 4 \times 10^{-3}$); NSGA-II and qEHVI, in contrast, are statistically indistinguishable ($p = 1.0$) and each one leads on different datasets.

\begin{table}[!ht]
    \centering
    \caption{Multi-objective \acrshort{nas}: mean hypervolume $\pm$ std over $N_{\mathrm{val}}=3$ seeds for the three native optimizers, with the gap to the true Pareto front $(\mathrm{HV}^{\star}-\mathrm{HV})$ in parentheses. The reference point is $R=(1.1,1.1)$; a higher hypervolume, and therefore a smaller gap, is better. The best value per dataset is in bold. The scalarized single-objective baseline is reported in Table~\ref{tab:study2_scalar_hv}.}
    \label{tab:study2_mo_hv}
    \resizebox{\textwidth}{!}{%
    \begin{tabular}{lccc}
        \toprule
        \textbf{Algorithm} & \textbf{CIFAR-10} & \textbf{CIFAR-100} & \textbf{ImageNet16-120} \\
        \midrule
        NSGA-II & $1.1321 \pm 0.0104$ \, ($8.1\!\times\!10^{-3}$) & $0.8940 \pm 0.0057$ \, ($6.9\!\times\!10^{-3}$) & $0.6053 \pm 0.0052$ \, ($5.3\!\times\!10^{-3}$) \\
        GDE3    & $\mathbf{1.1397 \pm 0.0002}$ \, ($4.1\!\times\!10^{-4}$) & $\mathbf{0.8993 \pm 0.0009}$ \, ($1.5\!\times\!10^{-3}$) & $\mathbf{0.6094 \pm 0.0012}$ \, ($1.3\!\times\!10^{-3}$) \\
        qEHVI   & $1.1362 \pm 0.0008$ \, ($3.9\!\times\!10^{-3}$) & $0.8942 \pm 0.0028$ \, ($6.6\!\times\!10^{-3}$) & $0.6035 \pm 0.0013$ \, ($7.2\!\times\!10^{-3}$) \\
        \midrule
        True Pareto front $\mathrm{HV}^{\star}$ & $1.1401$ & $0.9008$ & $0.6107$ \\
        \bottomrule
    \end{tabular}}
\end{table}

\paragraph{Coverage of the true front.}
The advantage of native search over scalarization is not uniform under hypervolume alone: on CIFAR-10 the dense low-cost fronts of the scalarized optimizers (Table~\ref{tab:study2_scalar_hv}) match NSGA-II. This is a limitation of the indicator rather than of the methods, since a tight cluster of cheap architectures can accumulate area without covering the front. Table~\ref{tab:study2_igd} reports IGD$^{+}$, which explicitly penalizes incomplete coverage. Under this indicator every native optimizer beats every scalarized one on every dataset, with GDE3 again first, so handling the two objectives natively recovers the high-accuracy region that scalarization systematically misses.

\begin{table}[!ht]
    \centering
    \caption{Multi-objective \acrshort{nas}: mean IGD$^{+}$ ($\times 10^{-3}$) $\pm$ std over $N_{\mathrm{val}}=3$ seeds to the exact Pareto front (lower is better), for the three native optimizers and the scalarized baseline of Part~B. The best value per dataset is in bold. Unlike hypervolume, IGD$^{+}$ explicitly penalizes incomplete coverage of the true front.}
    \label{tab:study2_igd}
    \setlength{\tabcolsep}{4pt}
    \resizebox{\textwidth}{!}{%
    \begin{tabular}{lcccccccc}
        \toprule
        & \multicolumn{3}{c}{\textbf{Native}} & \multicolumn{5}{c}{\textbf{Scalarized}} \\
        \cmidrule(lr){2-4} \cmidrule(lr){5-9}
        \textbf{Dataset} & \textbf{NSGA-II} & \textbf{GDE3} & \textbf{qEHVI} & \textbf{GA} & \textbf{SADE} & \textbf{GWO} & \textbf{AIW-PSO} & \textbf{BO} \\
        \midrule
        CIFAR-10       & $3.3 \pm 1.8$ & $\mathbf{1.0 \pm 0.3}$ & $2.5 \pm 0.4$ & $7.5 \pm 2.4$ & $5.2 \pm 0.3$ & $5.6 \pm 0.6$ & $6.3 \pm 0.3$ & $5.5 \pm 0.1$ \\
        CIFAR-100      & $3.7 \pm 1.1$ & $\mathbf{2.3 \pm 0.2}$ & $3.9 \pm 1.1$ & $8.8 \pm 0.1$ & $9.2 \pm 0.5$ & $8.7 \pm 0.1$ & $13.1 \pm 0.5$ & $9.8 \pm 1.6$ \\
        ImageNet16-120 & $4.3 \pm 1.3$ & $\mathbf{3.0 \pm 0.8}$ & $5.7 \pm 0.6$ & $10.7 \pm 0.2$ & $10.1 \pm 1.0$ & $9.6 \pm 1.2$ & $13.6 \pm 1.1$ & $9.0 \pm 0.6$ \\
        \bottomrule
    \end{tabular}}
\end{table}

\paragraph{Shape of the obtained fronts.}
Figure~\ref{fig:study2_mo_fronts} shows the Pareto front of each native optimizer in a separate panel, overlaid on the true front. NSGA-II and GDE3 follow the true front across its full extent, while qEHVI tracks it everywhere except at the most accurate, high-FLOPs end. Figure~\ref{fig:study2_mo_conv} reports the hypervolume as a function of the number of evaluations: all three methods converge close to the true-front level, with GDE3 both the highest and the most stable.

\begin{figure}[ht]
    \centering
    \includegraphics[width=\textwidth]{Chapters/img/nas_mo_fronts_native.pdf}
    \caption{Multi-objective \acrshort{nas}: Pareto front of each native optimizer shown in a separate panel and overlaid on the exact front (dashed). Each panel contains the three datasets, which occupy distinct accuracy bands.}
    \label{fig:study2_mo_fronts}
\end{figure}

\begin{figure}[ht]
    \centering
    \includegraphics[width=\textwidth]{Chapters/img/nas_mo_hv_convergence.pdf}
    \caption{Multi-objective \acrshort{nas}: hypervolume versus the number of evaluations, per dataset, for the three native optimizers (mean $\pm$ std over seeds). The dashed line is the true-front hypervolume; higher is better.}
    \label{fig:study2_mo_conv}
\end{figure}

\paragraph{The accuracy--cost trade-off in practice.}
Table~\ref{tab:study2_mo_tradeoff} makes the trade-off concrete along the best front obtained. On CIFAR-10, conceding about one accuracy point, from $0.9437$ to $0.9336$, already cuts the cost to $31\%$ of the most accurate model; a $0.91$-accuracy architecture needs only $13\%$ of its FLOPs, and $0.87$ needs $5\%$. The same pattern holds on CIFAR-100 and ImageNet16-120. The last column shows that NSGA-II and GDE3 recover the whole frontier, including the most accurate architecture, whereas qEHVI matches it everywhere except the top rows.

This is the practical payoff of formulating the problem with two objectives: the single-objective search of Part~A returns one architecture at the top of the accuracy axis, whereas Part~C returns the whole curve and lets the cost budget be decided after the search rather than before it.

\begin{table}[!ht]
    \centering
    \caption{Accuracy--cost trade-off along the best Pareto front obtained for multi-objective \acrshort{nas} (union of the NSGA-II, GDE3 and qEHVI fronts, aggregated over the three seeds), at representative points per dataset. \#FLOPs (\%) is relative to the most accurate architecture. The last column lists which methods recover each accuracy level (N: NSGA-II, G: GDE3, Q: qEHVI).}
    \label{tab:study2_mo_tradeoff}
    \setlength{\tabcolsep}{5pt}
    \begin{tabular}{lccccc}
        \toprule
        \textbf{Dataset} & \textbf{Top-1 acc.} & \textbf{\#FLOPs (M)} & \textbf{\#Params (M)} & \textbf{\#FLOPs (\%)} & \textbf{Reached by} \\
        \midrule
        \multirow{6}{*}{CIFAR-10} & 0.9437 & 153.3 & 1.073 & 100\% & NG \\
         & 0.9431 & 90.4 & 0.643 & 59\% & NG \\
         & 0.9391 & 78.6 & 0.559 & 51\% & NGQ \\
         & 0.9336 & 47.1 & 0.344 & 31\% & NGQ \\
         & 0.9126 & 19.6 & 0.157 & 13\% & NGQ \\
         & 0.8680 & 7.8  & 0.073 & 5\%  & NGQ \\
        \midrule
        \multirow{6}{*}{CIFAR-100} & 0.7351 & 184.7 & 1.294 & 100\% & NG \\
         & 0.7275 & 117.9 & 0.836 & 64\% & NGQ \\
         & 0.7163 & 55.0 & 0.406 & 30\% & NGQ \\
         & 0.7009 & 47.1 & 0.350 & 26\% & NGQ \\
         & 0.6708 & 15.7 & 0.135 & 8\%  & NGQ \\
         & 0.5887 & 7.8  & 0.079 & 4\%  & NGQ \\
        \midrule
        \multirow{6}{*}{ImageNet16-120} & 0.4731 & 30.5 & 0.866 & 100\% & NG \\
         & 0.4685 & 22.6 & 0.651 & 74\% & NGQ \\
         & 0.4540 & 13.7 & 0.407 & 45\% & NGQ \\
         & 0.4330 & 11.8 & 0.351 & 39\% & NGQ \\
         & 0.3833 & 3.9  & 0.136 & 13\% & NGQ \\
         & 0.2868 & 2.0  & 0.080 & 6\%  & NGQ \\
        \bottomrule
    \end{tabular}
\end{table}

\paragraph{Search cost.}
As with \acrshort{bo} in Part~A, qEHVI carries the overhead of its Gaussian-process surrogate: its runs took on the order of $2700$ seconds, against under $15$ seconds for NSGA-II and GDE3 (Table~\ref{tab:study2_time}), more than two orders of magnitude more, whereas the two evolutionary methods add negligible cost per evaluation.

\subsection{Discussion}
\label{subsec:study2_discussion}

The three parts of Study~2 expose complementary behaviors.

In the accuracy-only formulation, benchmark difficulty again drives the discrimination among optimizers. On CIFAR-10 and CIFAR-100 almost every method reaches the optimum; only ImageNet16-120 separates them, with \acrshort{bo} reaching the bound on every seed, \acrshort{sade} and \acrshort{gwo} approaching it within $10^{-3}$, \acrshort{ga} converging prematurely, and \acrshort{aiwpso} exploring widely but exploiting poorly. The architectural-diversity counts expose the mechanism behind each behavior, and the wall-clock measurements expose the price: \acrshort{bo} wins on quality but is about $450\times$ slower than any evolutionary method, because on a tabular benchmark the surrogate refitting is the entire cost of the search.

When the same optimizers are made cost-aware through scalarization, they trace nearly identical fronts that stay biased toward the low-cost region and never reach the high-accuracy architectures. The choice of optimizer becomes almost irrelevant compared with the choice of $\alpha$.

Handling the two objectives natively removes this bias. On hypervolume, GDE3 is significantly the best method (Wilcoxon, $p \approx 4 \times 10^{-3}$), while NSGA-II and qEHVI are indistinguishable from each other and, on CIFAR-10, even from the scalarized fronts. The coverage-sensitive IGD$^{+}$ indicator resolves this ambiguity: every native optimizer dominates every scalarized one on all datasets. GDE3 is the strongest method overall, pairing the smallest gap with the lowest variance across seeds and a per-run cost of a few seconds. Notably, it outperforms NSGA-II although both apply the same crowding-distance truncation, which points to the differential-evolution variation operator, rather than diversity preservation alone, as the differentiator.

Finally, the algorithm ranking established in Study~1 is not preserved here. \acrshort{gwo} is the strongest optimizer on the categorical space of Study~1, where the budget is small and each evaluation is expensive, yet on NAS-Bench-201 it is \acrshort{bo} that dominates the accuracy-only formulation, and on CIFAR-10 \acrshort{gwo} and \acrshort{ga} even swap the best and worst positions between the two studies. Two factors plausibly drive the reversal: the \acrshort{nas} budget is nearly five times larger, which favors the model-based exploitation of \acrshort{bo}, and the tabular objective is noise-free, which rewards a method that never re-evaluates the same point.

\clearpage
\section{Study 3 --- Hybrid Metaheuristic and Gradient-Based Weight Training}
\label{sec:study_3}

Studies~1 and~2 apply optimization to the \emph{configuration} of a network, its hyperparameters and its topology, while the weights themselves are always trained by gradient descent. The third study targets the remaining role of evolutionary methods inside the learning phase: the optimization of the weights.

\subsection{Motivation and Idea}
\label{subsec:study3_idea}

Gradient-based training is efficient at descending into a basin of attraction, but it is a local procedure: the solution it reaches depends on the initialization, and it offers no mechanism to escape a poor region of a non-convex, multimodal loss landscape. Population-based metaheuristics have the complementary profile: they explore the landscape globally and are insensitive to differentiability, but they converge slowly and refine poorly once a promising region has been located.

The idea of this study is to combine the two in sequence rather than to oppose them. A population-based optimizer is first used to search the weight space globally and bring the network into a promising region, an \emph{optimal zone}. The resulting weights are then handed over to a gradient-based optimizer, which performs the fine-grained refinement inside that region and drives the model to a local optimum. The metaheuristic supplies the exploration, the gradient method supplies the exploitation, and the point at which control is transferred becomes a parameter of the scheme.

\subsection{Planned Evaluation}
\label{subsec:study3_plan}

The evaluation is intended to reuse the infrastructure of Study~1: the same \acrshort{cnn} template, the same datasets, and the same convention of repeating every run over independent seeds and reporting distributions rather than single runs. The baseline is the same architecture trained purely by gradient descent from a standard initialization, under a training budget matched to the total budget of the hybrid scheme, so that any difference reflects the search strategy and not the amount of computation spent.

The main questions to be answered are whether the hybrid scheme reaches a better optimum than gradient descent alone under a matched budget, how sensitive the outcome is to the moment at which control is transferred from the metaheuristic to the gradient method, and whether any advantage survives as the number of weights grows.

\subsection{Status}
\label{subsec:study3_status}

This study is still under definition. The encoding of the weight vector, the choice of which layers are exposed to the metaheuristic, the switching criterion between the two phases, and the budget accounting that makes the comparison against a pure gradient baseline fair are open points, and the results are not yet available. The section is therefore presented as a plan; it will be completed once the experimental design is settled.


\section{Summary of the Experimental Evaluation}
\label{sec:experimental_summary}

Read together, the studies support four observations.

\paragraph{Benchmark difficulty determines what a comparison can show.}
MNIST in Study~1 and the two CIFAR datasets in Study~2 saturate against a ceiling and compress every algorithm into a near-indistinguishable band, in which secondary signals such as the population mean and the seed-to-seed dispersion carry more information than the headline result. CIFAR-10 in Study~1 and ImageNet16-120 in Study~2 expose clear and stable differences. An optimizer comparison is therefore only as informative as its hardest benchmark.

\paragraph{Algorithm rankings do not transfer across search spaces.}
\acrshort{gwo} is the strongest and most reliable optimizer on the categorical hyperparameter space of Study~1, where the budget is small and each evaluation requires training a network, yet \acrshort{bo} dominates the accuracy-only \acrshort{nas} formulation of Study~2, reaching the exhaustive optimum on every dataset and seed. Changing the formulation again, from single- to multi-objective, changes the answer once more, with GDE3 emerging as the best method. Optimizer selection should therefore account for the structure of the search space, the noise characteristics of the objective and the available budget, rather than rely on a fixed ranking.

\paragraph{Reliability and cost are first-class criteria.}
Averages under few repetitions can be misleading, as the \acrshort{ga} result on CIFAR-10 illustrates: the single best individual run coexists with the worst median and the largest variance among the population-based methods. Reporting distributions, medians and ranges guards against this. Cost deserves the same treatment: the Gaussian-process methods reach strong results but run more than two orders of magnitude slower per search than the evolutionary alternatives, an overhead that Study~1 could not isolate because it was hidden behind the cost of training the candidates, and that Study~2 makes explicit by removing training from the loop.

\paragraph{Self-adaptive formulations are the pragmatic evolutionary choice.}
\acrshort{sade} adapts its own control parameters and \acrshort{gwo} exposes none, so both ran without any configuration step while remaining competitive in both studies. This is a practical advantage in exactly the setting where automated search is needed, since a method that must itself be tuned reintroduces the problem it was meant to solve.

Study~3, once completed, will extend this reading to the third role of evolutionary methods inside the learning phase, by testing whether a global population-based search followed by gradient-based refinement improves on gradient descent alone under a matched budget.
